\subsection{序关系的定义}

设 \(\mathbb{R}\{x, y\}\) 为一个关系，其中 \(x\) 与 \(y\) 是不同的字母。 \(\mathbb{R}\) 被称为关于字母 \(x\) 与 \(y\) （或介于 \(x\) 与 \(y\) 之间）的一个序关系，如果

\[
(R \{x, y \} \text {   且   } R \{y, z \}) \Longrightarrow R \{x, z \},
\]

\[
(\mathrm{R} \{x, y \} \text {   且   } \mathrm{R} \{y, x \}) \Longrightarrow (x = y),
\]

\[
\mathrm{R} \left\{x, y \right\} \Longrightarrow (\mathrm{R} \left\{x, x \right\} \text { 且 } \mathrm{R} \left\{y, y \right\}).
\]

上述关系中的第一个表明，R关于字母x和y是传递的（第二章，§6，第1节）。

示例

\begin{enumerate}
\def\labelenumi{(\arabic{enumi})}
\item
  相等关系， \(x = y\) 是一个序关系。
\item
  关系 \(X \subset Y\) 是X与Y之间的一个序关系（第二章，§1，第2号，命题1和2以及公理A1），通常称为包含关系。
\item
  设 \(\mathbf{R}\{x, y\}\) 为 \(x\) 与 \(y\) 之间的一个序关系。关系 \(\mathbf{R}\{y, x\}\) 则是 \(x\) 与 \(y\) 之间的一个序关系，称为序关系 \(\mathbf{R}\{x, y\}\) 的相反关系。
\end{enumerate}

集合 E 上的一个序关系就是序关系 \(R\{x, y\}\) 关于两个不同的字母 x, y，使得关系 \(R\{x, x\}\) 等价于 \(x \in E\) （换言之，即满足 \(R\{x, y\}\) 在 E 上是自反的（第二章，§6，第 1 号））。则关系 \(R\{x, y\}\) 蕴含“ \(x \in E\) 并且 \(y \in E\) '' 以及关系 \((R\{x, y\}\) 并且 \(R\{y, x\})\) 等价于“ \(x \in E\) 并且 \(y \in E\) 并且 \(x = y\) ''.

示例

\begin{enumerate}
\def\labelenumi{(\arabic{enumi})}
\item
  相等关系和包含关系并不是集合上的序关系，因为这些关系 \(x = x\) 与 \(\mathbf{X} \subset \mathbf{X}\) 并非收集性的（第二章，§1，第7号）。
\item
  设 \(\mathbb{R}\{x, y\}\) 为 \(x\) 与 \(y\) 之间的一个序关系，并且设 \(\mathbb{E}\) 为一个集合，使得 \(x \in \mathbb{E}\) 蕴含 \(\mathbb{R}\{x, x\}\) (注意，空集满足此条件)。关系“ \(\mathbb{R}\{x, y\}\) 并且 \(x \in \mathbb{E}\) 并且 \(y \in \mathbb{E}\) '' 于是是 \(\mathbb{E}\) 上的一个序关系，这一点立即可以验证；它被称为由 \(\mathbb{R}\{x, y\}\) 在 \(\mathbb{E}\) 上（参见第4号）诱导的序关系。滥用语言时，短语“关系 \(S\{x, y\}\) 是 \(\mathbb{E}\) 的元素之间的序关系”常被用来代替“关系（ \(S\{x, y\}\) 并且 \(x \in \mathbb{E}\) 并且 \(y \in \mathbb{E}\) ）是 \(\mathbb{E}\) '' 上的一个序关系。例如，如果 \(A\) 是一个集合，关系“ \(X \subset Y\) 并且 \(X \subset A\) 并且 \(Y \subset A\) '' 是 \(A\) 的子集之间的一个序关系.
\item
  设E，F为集合。关系“g延拓f”是E的子集到F的映射之集合上的一个序关系（在f与g之间）。
\item
  在集合 \(\mathfrak{P}(\mathfrak{P}(E))\) （由一个集合 \(E\) 的子集集合所组成）中，设 \(\mathfrak{F}\) 为 \(E\) 的划分所成的集合（第二章，§4，第7号）。我们回顾一下，一个划分 \(\varpi\) 被称为比一个划分 \(\varpi'\) 更粗糙，若给定任意 \(Y \in \varpi'\) ，存在 \(X \in \varpi\) 使得 \(Y \subset X\) (第二章，§4，第6号)。对于每个划分 \(\varpi\) （ \(E\) 的划分），令 \(\tilde{\varpi}\) 为 \(\varpi\) 在 \(E\) 上（第二章，§ 6，第 2 号）定义的等价关系的图，即（互不相交的）集合 \(A \times A\) 的并集，其中 \(A\) 贯穿 \(\varpi\) 。关系“ \(\varpi\) 比 \(\varpi'\) 更粗”立即被看出等价于 \(\tilde{\varpi} \supset \tilde{\varpi}'\) ，因此是集合 \(\mathfrak{F}\) 上 \(\varpi\) 与 \(\varpi'\) 之间的一个序关系。
\end{enumerate}

集合 E 上的一个序是一种对应关系 \(\Gamma = (G, E, E)\) 以E为源，以E为目标，且该关系 \((x, y) \in G\) 是E上的一个序关系。滥用语言时，我们有时会将 \(\Gamma\) 的图G称为E上的一个序。如果 \(R\{x, y\}\) 是 E 上的一个序关系，则它具有一个图，该图是 E 上的一个序。

命题1. 一个介于 E 与 E 之间的对应关系 \(\Gamma\) 是 E 上的一个序，当且仅当其图 G 满足以下条件：

\begin{enumerate}
\def\labelenumi{(\alph{enumi})}
\item
  \(\mathbf{G} \circ \mathbf{G} = \mathbf{G}\) .
\item
  集合 \(\mathbf{G} \cap \overline{\mathbf{G}}\) 是 \(\Delta\) ，即 \(\mathbf{E} \times \mathbf{E}\) 的对角线.
\end{enumerate}

对于关系 \(((x, y) \in \mathbf{G}\) 并且 \((y, z) \in \mathbf{G}) \Longrightarrow ((x, z) \in \mathbf{G})\) 可以写成 \(\mathbf{G} \circ \mathbf{G} \subset \mathbf{G}\) ，而关系

\((x, y) \in G\) 并且 \((y, x) \in G) \iff (x = y\) 并且 \(x \in E\) 并且 \(y \in E)\)

可以写成 \(\mathbf{G} \cap \overline{\mathbf{G}} = \Delta\) 。由 \(\mathbf{G} \cap \overline{\mathbf{G}} = \Delta\) 我们于是推出 \(\Delta \subset \mathbf{G}\) ，由此 \(\mathbf{G} = \Delta \circ \mathbf{G} \subset \mathbf{G} \circ \mathbf{G}\) ；又因为 \(\mathbf{G} \circ \mathbf{G} \subset \mathbf{G}\) ，我们有

\[
\mathbf {G} \circ \mathbf {G} = \mathbf {G}.
\]

